% META: {"id": "TEXP-MAT-ME-002", "chapitre": "TEXP-MATRICES-MARKOV", "type_objet": "methode", "capacites_codes": ["C2"], "methodes": ["M2"], "mode_creation": "ex_nihilo", "sources_inspiration": [], "status": "needs_review"}
% BEGIN-VERIFY
% from sympy import Matrix
% # 2 usines x 3 produits ; production journaliere
% P = Matrix([[10, 5, 8], [4, 12, 6]])
% # totaux sur 5 jours
% assert 5*P == Matrix([[50, 25, 40], [20, 60, 30]])
% # couts unitaires des 3 produits -> cout par usine = P * couts
% couts = Matrix([2, 3, 1])
% assert P*couts == Matrix([43, 50])
% END-VERIFY
\begin{fichemethode}{M2}{Modeliser une situation par une matrice}

\quandUtiliser{%
  Utiliser cette methode lorsqu'une situation a DEUX entrees (usines/produits,
  liaisons d'un reseau, systeme lineaire...) doit etre codee par une matrice
  avant tout calcul.
}

\pasApas{%
  \begin{enumerate}
    \item Choisir la signification des LIGNES et celle des COLONNES, et
      l'annoncer explicitement (ex. : lignes = usines, colonnes = produits).
    \item Remplir la matrice : le coefficient $(i\,;\,j)$ est la donnee croisee
      de la ligne $i$ et de la colonne $j$ (production, cout, effectif...).
    \item Traduire les operations de l'enonce : cumul dans le temps $\to$
      multiplication par un scalaire ou somme ; bilan croise (couts,
      totaux) $\to$ produit par un vecteur colonne ou une autre matrice.
    \item Controler la COHERENCE des dimensions avant chaque produit : c'est
      le meilleur detecteur d'erreur de modelisation.
    \item Interpreter chaque coefficient du resultat dans le contexte (unites
      comprises).
  \end{enumerate}
}

\exempleRedige{%
  Deux usines fabriquent trois produits ; les productions journalieres sont :
  usine 1 : 10, 5, 8 ; usine 2 : 4, 12, 6. Les couts unitaires des produits
  sont 2, 3, 1 (en euros). Modeliser et calculer le cout journalier par usine.

  \medskip
  \commentaireMarge{Lignes $=$ usines, colonnes $=$ produits.}%
  $P = \begin{pmatrix} 10 & 5 & 8\\ 4 & 12 & 6 \end{pmatrix}$ et le vecteur
  des couts $C = \begin{pmatrix} 2\\ 3\\ 1 \end{pmatrix}$.

  \medskip
  \commentaireMarge{Dimensions : $(2\times 3)(3 \times 1) = 2 \times 1$.}%
  $PC = \begin{pmatrix} 10 \times 2 + 5 \times 3 + 8 \times 1\\
  4 \times 2 + 12 \times 3 + 6 \times 1 \end{pmatrix}
  = \begin{pmatrix} 43\\ 50 \end{pmatrix}$ : 43 euros par jour pour l'usine 1,
  50 pour l'usine 2.
}

\pieges{%
  \begin{itemize}
    \item \textbf{Piege 1.} Ne pas annoncer la convention lignes/colonnes : la
      matrice transposee code une autre lecture et fausse les produits.
    \item \textbf{Piege 2.} Multiplier dans le mauvais sens ($CP$ au lieu de
      $PC$) : les dimensions ne s'enchainent plus.
    \item \textbf{Piege 3.} Additionner des matrices codant des grandeurs
      d'unites differentes.
  \end{itemize}
}

\verifier{%
  Recalculer un coefficient du resultat directement depuis l'enonce (43 euros
  pour l'usine 1), et verifier que dimensions et unites sont coherentes.
}

\sEntrainer{\refExos{M2}}

\end{fichemethode}
